\section{Step 2: Definition of the indicators}
\label{sec:brainstorming_indicators}
The initial phase of this \gls{mcda} study focuses on indicator selection and hierarchical structuring. The following sections outline the methodological process employed to develop the final indicator set. It should be noted that this study primarily evaluates the feasibility of the proposed approach rather than generating definitive results. Accordingly, this study serves to identify which indicators or areas may require further detailed development in future iterations, as the weighting process will naturally reveal priority areas. Indicators receiving minimal weights may not justify further refinement, while those with substantial weights could benefit from more granular analysis. With this context established, the subsequent sections present the process that resulted in the final indicator set used in this study.


Following the bottom-up approach outlined in Section \ref{sec:mavt_presentation}, the initial analytical step involved identifying all potential differentiating factors among the six reactor alternatives. The process began with a series of personal brainstorming sessions by the author, followed by a preliminary selection phase. The resulting comprehensive list of 35 potential indicators\footnote{The complete list is available as a CSV file in the repository, with an example Python script provided for navigation (see Section \ref{sec:appendix_code} for details)} was then extensively reviewed with the support of three domain experts to evaluate:
\begin{itemize}
    \item Completeness (identifying missing aspects)
    \item Redundancy (eliminating repetitive measures)
    \item Relevance (removing irrelevant indicators)
    \item Granularity (balancing overly specific or aggregate measures)
\end{itemize}

Further discussion with several experts from different backgrounds was helpful to identify other possible aspects, as well as to collect feedback that helped improve the justifications. On this regard it is really important that the analyst takes into serious consideration every suggestion provided by others, and, whenever these have to be discarded, it is important to report it and properly justify the omission to ensure full transparency.


As anticipated, the extensive indicator list required systematic screening for several reasons:
\begin{itemize}
    \item Double counting: The same aspects might be represented by multiple indicators, leading to biased weights
    \item Scope: Some aspects might be valid but not relevant for the point of view of the study, in this case, European governmental agencies
    \item Data availability: This wouldn't be the case for a field application but for the purpose of this thesis obtaining data outside the public domain was deemed unfeasible.
\end{itemize}

Table \ref{tab:screening-indicators} provides examples of excluded indicators. The most significant exclusions are justified in Section \ref{sec:indicators_removed} while possible improvements for future studies are documented in Section \ref{sec:results}.


\begin{table}[!hb]
\centering
\begin{tabularx}{0.8\textwidth}{Xccccc}
\toprule
\textbf{Indicator} & \textbf{F1} & \textbf{F2} & \textbf{F3} & \textbf{Decision} \\
\midrule
Load Following Capabilities & \cmark & \cmark & \cmark & \textbf{Selected} \\

Mining waste & \cmark & \cmark & \xmark & Rejected \\

Land Usage & \xmark & \xmark & \cmark & Rejected \\

Core Damage Frequency & \xmark & \cmark & \cmark & Rejected \\

Spent Nuclear Fuel & \cmark & \cmark & \cmark & \textbf{Selected} \\

Local Economy Turnaround & \xmark & \xmark & \cmark & Rejected \\

Capital Investment Budgeted & \cmark & \cmark & \cmark & \textbf{Selected} \\

Supplier Availability  & \cmark & \cmark & \cmark & \textbf{Selected} \\
\bottomrule
\end{tabularx}
\caption{Example of the selection process of indicators. F1: High relevance to the case study, F2: Data is available, F3: Not redundant nor highly correlated with other indicators}
\label{tab:screening-indicators}
\end{table}

Following this systematic screening process, twelve indicators remained and were subsequently aggregated into higher-level objectives, presented in Table \ref{tab:final_indicators}. In this bottom-up approach, each group collects indicators that share a common high-level point of view. The subsequent sections describe each retained indicator, covering its data sources, conceptual explanations, and the methodological rationale for its inclusion.

\begin{table}[!hb]
    \centering
    \begin{tabularx}{\textwidth}{lllX}
        \toprule
        \multicolumn{4}{c}{\textbf{Technical Indicators}} \\
        \midrule
        T1 & Thermal Efficiency & [\%] & Indicator of the thermal hydraulics efficiency \\
        \cmidrule(lr){1-4}
        T2 & Discharge Burnup & [GWd/tonHM] & Indicator of the fuel utilization efficiency and neutron economy \\
        \cmidrule(lr){1-4}
        T3 & Net Power Output & [MWe] & How much does the output of one unit fit the grid demand \\
        \cmidrule(lr){1-4}
        T4 & Load Following Capabilities & [-] & How much do the load following capabilities satisfy the grid demand \\
        \midrule
        \multicolumn{4}{c}{\textbf{Feasibility Indicators}} \\
        \midrule
        F1 & Design Maturity & [Score] & Technological readiness of the design \\
        \cmidrule(lr){1-4}
        F2 & Licensing Status & [Score] & How close to having a licence is the reactor? \\
        \cmidrule(lr){1-4}
        F3 & Supplier Availability & [Score] & Is the supplier established in the country? Even partially, for instance has offices \\
        \midrule
        \multicolumn{4}{c}{\textbf{Economic Indicators}} \\
        \midrule
        E1 & Capital Investment Budgeted & [2024EUR/kWe] & Overnight, budgeted construction costs (estimates include decommissioning etc.) \\
        \cmidrule(lr){1-4}
        E2 & Design Complexity & [Score] & From which delays in the manufacturing processes might arise \\
        \cmidrule(lr){1-4}
        E3 & Construction Complexity & [Score] & From which delays in the building processes might arise \\
        \cmidrule(lr){1-4}
        E4 & Operative and Fuel Costs & [USD2020/MWh] & Ongoing costs \\
        \midrule
        \multicolumn{4}{c}{\textbf{Social Indicators}} \\
        \midrule
        P1 & Electricity Market Benefit & [2024EUR/MWh] & Potential reduction of cost in the Electricity Market by marginal price logic, by constructing one unit \\
        \cmidrule(lr){1-4}
        P2 & GHGe Benefit & [tonCO2eq] & Potential Yearly Reduction on GHG Emissions by installing one unit \\
        \cmidrule(lr){1-4}
        P3 & Perceived Safety & [Score] & Based on public opinion, how much each reactor sounds safe? \\
        \cmidrule(lr){1-4}
        P4 & Nuclear Waste & [tonSNF] & Estimation of amount of spent fuel produced \\
        \bottomrule
    \end{tabularx}
    \caption{Final selection of indicators, categorized and coded for evaluation.}
    \label{tab:final_indicators}
\end{table}
